\documentclass[12pt]{article}
 \usepackage[german]{babel}
 \usepackage[utf8]{inputenc}
 \usepackage{graphicx}
 \author{Denis Herdt, Almin Causevic}
 \title{ROS auf dem Rapberry Pi}
 \setlength{\parindent}{0pt}                   % Einrueckung 1. Zeile eines Absatzes
 \setlength{\parskip}{5pt plus 2pt minus 1pt}  % Abstand zwischen Absaetzen
 \frenchspacing
 \sloppy
 
 
\begin{document}
%4-20 Seiten
\begin{figure}[h]

\includegraphics[width=4cm]{hs-logo.jpg}
\end{figure}
Fakultät für Elektrotechnik und Informatik

\vspace{3cm}

\begin{center}

{\bf \huge ROS auf dem Raspberry Pi}
\vspace{4cm}

14 November, 2013
\vspace{1cm}

Systemadministration Projekt in Angewandter Informatik \\
von Denis Herdt und Almin Causevic

\end{center}

\pagebreak

\tableofcontents

\pagebreak

\section{Einleitung}
1-2 Seiten
\subsection{Motivation}

Wir haben uns für dieses Thema entschieden, weil wir in Kombination mit Hard- Software arbeiten wollten und mitbekommen haben, wie umständlich z.B. der Roboter ''Volksbot'' im Robotiklabor der Hochschule mit einem zusätzlichen Laptop gesteuert werden muss. Deshalb haben wir uns entschlossen, dieses Problem mithilfe des Raspberry Pi zu lösen. Die Idee stieß auch auf großes Interesse bei den Labormitarbeitern und bringt viel praktischen Nutzen.

Unser Ziel möchten wir mithilfe des Betriebssystems ROS (Robot Operating System) realisieren.
Für ROS haben wir uns entschieden, weil es zur Zeit die aktuellste und inovativste Methode ist, ein flexibles Netzwerk von Komponenten für z.B. die Robotersteuerung zu realisieren. Es wird Weltweit eingesetzt und könnte auch in Zukunft sehr nützlich für die Forschung werden. %für wen nützlich(uns,welt,niemand??)

Durch die Kombination mit der Mächtigkeit und Flexibilität von ROS und der geringen Größe und dem niedrigen Energieverbrauch des Pi eröffnen sich außerdem neue und effizientere Einsatzmöglichkeiten.

Sehr gut finden wir auch die Tatsache, dass wir mit Linux arbeiten können und dieses Projekt viel Konfiguration des Betriebssystems voraussetzt, da wir auch gerne mehr Praxisanwendung mit Linux haben möchten.


%Roboter Lab Anwendung von ROS und Möglichkeit, Roboter zu nutzen

%Interesse der Lab-Crew an Umsetzung ROS mit Raspberry Pi

%Kombination Hard-/Software

%Einarbeitung in Kommunikations Framework ROS

%Mit Hardware in Berührung zu kommen

%Konfiguration mit Linux

%Nodes nun mit Hardware (Pi) realisierbar

\subsection{Zielstzung}

Recherche über Linux Systeme auf Raspberry Pi
Passendes Linux System auf Raspberry Pi aufsetzen					
Recherche über ROS									
geeigneten Roboter wählen (Recherche Verfügbarkeit, Komponenten etc.)	
eventuell Komponenten beschaffen (Wlan Stick, etc.)				
ROS auf Raspberry Pi aufsetzen
Wlan-Netzwerk zwischen 2 ROS-kompatiblen Systemen herstellen
Netzwerk mithilfe von Testausgaben überprüfen \\

{\bf optional}
Roboter durch einen Raspberry Pi mithilfe von ROS steuern
Bild einer Webcam durch ROS auf einen Bildschirm streamen

\subsection{Eigene Leistung}

anstatt Linux PC wird Pi benutzt

\subsection{Aufbau der Arbeit}

Recherche über Linux Systeme auf Raspberry Pi
Passendes Linux System auf Raspberry Pi aufsetzen					
Recherche über ROS									
geeigneten Roboter wählen (Recherche Verfügbarkeit, Komponenten etc.)	
eventuell Komponenten beschaffen (Wlan Stick, etc.)				
ROS auf Raspberry Pi aufsetzen		
Wlan-Netzwerk zwischen 2 ROS-kompatiblen Systemen herstellen
Netzwerk mithilfe von Testausgaben überprüfen \\

{\bf optional}
Roboter durch einen Raspberry Pi mithilfe von ROS steuern
Bild einer Webcam durch ROS auf einen Bildschirm streamen

\section{Grundlagen}
%6-8 Seiten

\subsection{Raspberry Pi Hardwarekomponenten}

Der Raspberry Pi ist ein kreditkartengroßer Einplatinencomputer, der von der Raspberry Pi Foundation entwickelt wurde.
Wir benutzen für dieses Projekt das leistungsstärkere Modell B.

Technische Details:

\begin{itemize}
\item Preis: ca. 35 Euro
\item Prozessor: ARM1176JZF-S (700 MHz)
\item Broadcom VideoCore IV
\item SDRAM: 512 MB
\item Bis zu 16 GPIO-Pins
\item USB-Anschlüsse: 2
\item FBAS, HDMI
\item 3,5-mm-Klinkenstecker (analog), HDMI (digital)
\item Kartenleser für SD (SDHC und SDXC)/MMC/SDIO
\item 10/100-MBit-Ethernet-Controller 
\item 5 V, 700 mA (3,5 Watt)
\item 5-V-Micro-USB-Anschluss (Micro-B), alternativ 4 x AA-Batterien
\end{itemize}

Für den Anschluss an einen Bildschirm benutzen den HDMI Ausgang.
Diesen benutzen wir für Kontrollausgaben und Tests.

Als Stromquelle steht ein Akkupack ......(Firma,Modell) zum Einsatz. 
Hierbei ist wichtig, dass der Raspberry mind. 700 mA, besser 1 A zur Stromversorgung bekommt. 
Bei niedrigerer Amperzahl arbeitet der PC oft nicht zuverlässig.

Ein schneller und großer RAM Speicher ist für unsere Zwecke wichtig, da viele Signale, teils sogar synchron, verarbEin 17-Zoll Monitor wird für die Bildschirmausgabe des Raspberry über HDMI zu DVI benutzteitet werden müssen.

Für die Datenübertragung über das Netz benutzen wir einen Standard 300Mbit/s Wlan-Stick.

Der Raspberry Pi arbeitet mit einer ARM-Prozessorarchitektur.
Diese Architektur wird gerne für embedded Systems, wie PDAs oder Router, eingesetzt.

Sie ist auch auf jedem Smartphone zu finden, da sie den Vorteil einer sehr geringen Leistungsaufnahme bietet. 

Erwähnenswert ist die Architektur deshalb, weil wir im Laufe des Projektes Schwierigkeiten hatten, auf die wir später eingehen werden
Wir benutzen das auf Linux basierende Betriebssystem Raspbian.
Dabei handelt sich um ein für Raspberry Pi optimiertes open-source Debian-System.
Es enthält viele für die ARM-Architektur vorkompilierte Pakete (über 35.000), dazu auch Features wie etwa eine GUI.

Das System braucht 3GB Speicherbedarf unserer 16GB großen SD-Karte.
Der zusätzliche Speicherplatz wäre nötig, falls man vorhat, Log-Dateien und Ähnliches direkt auf dem Raspberry Pi zu speichern. 
In unserem Fall übernehmen das jedoch die leistungsstärkeren PC's über das ROS Netzwerk.

GPIO (General Purpose Input/Output) ist ein weiteres interessantes Feature des Raspberry Pi.
Es gibt uns die Möglichkeit, jegliche Hardware-Funktionalität anzusteuern.
Beispielsweise können LED-Leuchten oder der Start-Knopf der Kaffeemaschine damit über ROS gesteuert werden.
Wir konnten uns leider jedoch zeitlich bedingt nicht mehr mit dieser Thematik beschäftigen, es bietet aber Anreiz für noch mehr Ansätze und Umsetzungen mithilfe des Raspberry Pi und ROS.

%%%%%
%GPIO-Pins
%%%%%

%Wlan Übertragung für Raspberry über Stick sehr leistungsfressend !

32-Bit-Chip-Design
Aufgrund ihrer geringen Leistungsaufnahme kommen ARM-Prozessoren vor allem in eingebetteten Systemen, wie Mobiltelefonen, PDAs und Routern zum Einsatz.
Besondere Bedeutung hat die Architektur bei Smartphones
Zukunft: ARMv8 64-Bit 2014, iPhone 5s soll darauf basierenden Chip benutzen.
Unsere PC's basieren auf einer x64 (64Bit) oder x84 (32Bit) Architektur.
%%%%%
% ARM-Architektur 
% x_86 Architektur
% Unterschiede
%%%%%


\subsection{Smartphone}
%App-Steffen

\subsection{Volksbot}

\subsection{ROS}
%catkin, Netwerk



\section{Problem}

Hauptproblem??

\subsection{Verwandte Arbeiten}

Steffen, Marc
ROS-Tut

\section{Anforderungen}

Datenpakete übers ROS Netzwerk verschicken kann
Beliebige Datentypen
Beliebige Linux-Systeme
Überprüfung davon Log-Dateien, Kontrollausgaben etc..
robuste und zuverlässige Kommunikation möglich
modulares und flexibles Netzwerk
Pakete(repos) sollen leicht in Projekte eingebunden werden

optional:
Roboter soll bei Steuerung in richtige Richtungen fahren
Roboter soll sofort auf Richtungsanweisungen reagieren
Robotersteuerung Sensibilität sollte einstellbar sein
Steuerung über Wlan
Raspberry Pi auf Roboter ohne Kabelzgewirr
Node des Raspberry starte automatisch mit Raspbian
Leichte Bedienbarkeit des Roboters

optional:
Streaming-Bild ruckelfrei bearbeitet werden
soll unkonvertiert in Echtzeit übertragen werden
soll an externem Bildschirm ausgegeben werden

\section{Lösungsvorschläge}

Netzwerk wird mit ROS-Master aufgesetzt, über topics und nodes
verschicken und empfangen von Paketen 

beliebige Datentypen durch msg und srv Funktion realisierbar

ROS unterstüzt beliebige Linux Systeme..

Horchen an Topics mit echo + Ausgabe im Code mithilfe ROS-Stream Fkt.

Begrenzung der Anzahl der Datenübertragung
Verschickte und erwartete Datentypen müssen stimmen
Richtige Konfiguration des Masters und Clients

durch ROS und catkin gegeben 

durch catkin gegeben

opt:
abhängig vom richtigen Code der Motorsteuerung

abhängig von Raspberry Pi Hardware und effizientem Code und
guter Hardware

abhängig vom Code

Hardware nötig Wlan-Stick, Akku-Pack 

sollte in Linux beim booten konfiguriert werden

gegeben durch gute Smartphone App

opt:
abhängig von Raspberry Pi Hardware und Übertragungs-Codierung

abhängig vom Codex Hardware Komponenten

nötige Hardware und Netzwerk legen

\section{Bewertung der Lösungen}

alles gut

\section{Implementation}

ROS Netzwerk (topics, nodes, catkin...)
Linux SD Karte Pi
Datenübertragung an Volksbot
Konfiguration Motorsteuerung-Paket und Compilieren
Kontrollausgaben 
Netzwerkkomponenten zusammen arbeiten lassen (Hard/-Software)
Wlan -> Teilproblem
Kamera -> Teilproblem Zeit


\section{Fazit}
1 Seite
\subsection{Zusammenfassung}

sehr knapper Zeitrahmen
hohe Anforderungen für diesen Zeitraum
interessantes Projekt
anspruchsvolles Projekt
viel über Linux und Ros und Hardware und Netzwerken gelernt
Spaß

\section{Ausblick}

Benutzung des Pi's als Hardware Nodes Lab
Benutzung des Pi's für kleine Roboterprojekte Lab
Immer mehr binary packages für ROS Pi verfügbar -> mehr Möglichkeiten
z.B.: GPIO Interface -> Knöpfe Kaffemaschine, Relays, alle möglichen Signale
auf Hardwareebene

\section{Eigene Leistung}

anstatt Linux PC wird Pi benutzt

\section{Quellen}
\begin{verbatim}

http://de.wikipedia.org/wiki/Raspberry_Pi
http://www.amazon.de/Raspberry-Pi-RBCA000-Mainboard-1176JZF-S/dp/B008PT4GGC
http://www.raspbian.org/
http://www.softwareok.de/?seite=faq-System-Allgemein&faq=13
http://de.wikipedia.org/wiki/ARM-Architektur
http://www.rn-wissen.de/index.php/Raspberry_PI:_GPIO

\end{verbatim}

\end{document}